% Shared exercise chunk: l4-sharing-costs
\begin{exercisechunk}
\item \textbf{Task costs at a given learning rate.\suppmark}
\label[exercise]{ex:sharing-costs}
Prove \cref{prop:sharing-full-cost} using the following construction.
At each horizon \(T\), partition
the first \(T\) positions into \(\lceil d(T)\rceil\) nonempty groups.
Let the \(U_t\) be independent fair bits and \(V_t=0\) under \(P\).
Use terminal cost \(V_T\) and counting cost \(\sum_{t=1}^T V_t\).
Write \(P_T\) and \(P_T^U\) for the laws of \((U,V)_{1:T}\) and
\(U_{1:T}\), respectively.
If \(\widehat P^U_T\) is the pooled
smoothed estimate of the law of \(U_{1:T}\), set
\[
 r_T=\E\!\left[\KL(P^U_T,\widehat P^U_T)\right]
\]
and, for any law $R$ supported on $v_{1:T}\ne0^T$, define
\[
 \widehat Q_R
 =e^{-r_T}(\widehat P^U_T\otimes\delta_{0^T})
  +(1-e^{-r_T})R.
\]
\begin{enumerate}[(a),beginpenalty=10000]
\item Show that \(r_T\) is determined by \(n,T\) and the chosen
partition, so the mixture weight can be used by the estimator.
\item Compute \(\E[\KL(P_T,\widehat Q_R)]\) and show that it has
order \(d(T)/n\), independently of \(R\).
\item Determine the full attainable intervals for the expected terminal
and counting costs as \(R\) varies. Identify laws attaining the endpoints.
\end{enumerate}

\end{exercisechunk}
